\chapter{התוכנית הפיננסית}\label{ch:financials}

הפרק מציג את המודל הפיננסי של RoofRoot לחמש שנים, \RRroofsFirst{} גגות ב-2027 ו-\RRroofsFinal{} בסוף 2031. כל המספרים בו נגזרים ממודל אחד, שנשען על הנחות שמוצגות בטבלה~\ref{tab:assumptions}. המודל הוא אוסף של הנחות ולא תחזית שמתחייבים אליה: אם ההנחות משתנות, התוצאות משתנות איתן, ובסעיף הרגישות נראה בכמה.

\section{הנחות היסוד}

ההנחות מחולקות לשלוש קבוצות. הראשונה היא כלכלת גג בודד: שטח גידול של \RRarea{} מ"ר, תפוקה של \RRyield{} ק"ג למ"ר בשנה, כלומר \RRroofKg{} ק"ג לגג, במחיר ממוצע של \RRprice~\shekel{} לק"ג. התפוקה שמרנית ביחס למדווח בספרות על מערכות \en{NFT} \cite{barbosa2015}. ההכנסה השנתית של גג בוגר היא \RRroofRevenue{} אלף \shekel{}, והשקעת ההקמה היא \RRcapexRoof{} אלף \shekel{}.

הקבוצה השנייה היא עלויות משתנות: דמי גג בשיעור \RRrentShare{} אחוזים מההכנסה שמועברים לוועד הבית, ועלות תפעול קבועה לכל גג. הקבוצה השלישית היא הנחות פיננסיות: שיעור היוון של \RRdiscount{} אחוזים, מכפיל יציאה של \RRexitMultiple{} על \en{EBITDA}, מס חברות ופחת בקו ישר.

\begin{table}[htbp]
\centering\small
% AUTO-GENERATED by figures/python/make_figures.py — do not edit by hand.
\rrstripes
\begin{tabularx}{\linewidth}{Ycc}
\rrtablehead \rrth{הנחה} & \rrth{ערך} & \rrth{יחידה} \\
שטח גידול נטו לגג & 250 & מ"ר \\
יבול שנתי & 30 & ק"ג למ"ר \\
מחיר ממוצע משוקלל & 24 & \shekel\ לק"ג \\
ניצולת בשנה הראשונה של גג & 50\% & — \\
דמי גג לוועד הבית & 8\% & מההכנסות \\
עלות תפעול לגג & \en{72} & אלפי \shekel\ לשנה \\
השקעה הונית לגג & \en{180} & אלפי \shekel \\
רישיון \en{RootOS} & 1,500 & \shekel\ לחודש \\
שיעור היוון & 12\% & — \\
מכפיל יציאה & 8× & \en{EBITDA} \\
\end{tabularx}

\caption{הנחות היסוד של המודל. כל ערך בטבלה מוזן מקובץ הנתונים של התוכנית.}
\label{tab:assumptions}
\end{table}

גג שנפתח באמצע התוכנית אינו מגיע מיד לתפוקה מלאה, ולכן בשנת ההקמה מניחים שהוא פועל בחצי קיבולת. גג בוגר מניב תרומה של \RRroofContribution{} אלף \shekel{} בשנה אחרי דמי הגג ועלות התפעול, והחזר ההשקעה בו הוא \RRroofPayback{} שנים.

\section{תוצאות: דוח רווח והפסד}

טבלה~\ref{tab:pnl} מציגה את התוצאות לפי שנים. ההכנסות גדלות מ-\RRrevFirst{} מיליון \shekel{} ב-2027 ל-\RRrevFinal{} מיליון ב-2031. עיקר ההכנסה מגיע ממכירת תוצרת, וחל מ-2029 גם מרכיב של רישוי \en{RootOS} לחוות חיצוניות, שמגיע ל-\RRsaasFarmsFinal{} חוות בסוף התקופה.

\begin{table}[htbp]
\centering\small
% AUTO-GENERATED by figures/python/make_figures.py — do not edit by hand.
\rrstripes
\begin{tabularx}{\linewidth}{Yrrrrr}
\rrtablehead \rrth{\shekel\ אלפים} & \rrth{2027} & \rrth{2028} & \rrth{2029} & \rrth{2030} & \rrth{2031} \\
גגות פעילים (סוף שנה) & 4 & 12 & 28 & 50 & 80 \\
הכנסות תוצרת & \en{360} & \en{1,440} & \en{3,600} & \en{7,020} & \en{11,700} \\
הכנסות \en{SaaS} & 0 & 0 & \en{270} & \en{810} & \en{1,620} \\
\textbf{סך הכנסות} & \textbf{\en{360}} & \textbf{\en{1,440}} & \textbf{\en{3,870}} & \textbf{\en{7,830}} & \textbf{\en{13,320}} \\
דמי גג לוועד הבית & \en{-29} & \en{-115} & \en{-288} & \en{-562} & \en{-936} \\
עלות תפעול גגות & \en{-144} & \en{-576} & \en{-1,440} & \en{-2,808} & \en{-4,680} \\
\textbf{רווח גולמי} & \textbf{\en{187}} & \textbf{\en{749}} & \textbf{\en{2,142}} & \textbf{\en{4,460}} & \textbf{\en{7,704}} \\
מטה והנהלה & \en{-1,600} & \en{-2,200} & \en{-3,000} & \en{-3,800} & \en{-4,500} \\
\textbf{\en{EBITDA}} & \textbf{\en{-1,413}} & \textbf{\en{-1,451}} & \textbf{\en{-858}} & \textbf{\en{660}} & \textbf{\en{3,204}} \\
פחת & \en{-144} & \en{-432} & \en{-1,008} & \en{-1,800} & \en{-2,880} \\
מס & 0 & 0 & 0 & 0 & \en{-75} \\
\textbf{רווח נקי} & \textbf{\en{-1,557}} & \textbf{\en{-1,883}} & \textbf{\en{-1,866}} & \textbf{\en{-1,140}} & \textbf{\en{249}} \\
\end{tabularx}

\caption{דוח רווח והפסד ותזרים מצומצם, באלפי \shekel{}. המספרים השליליים הם הוצאות.}
\label{tab:pnl}
\end{table}

שולי הרווח הגולמי בשנה האחרונה הם \RRgrossMarginFinal{} אחוזים, ושולי \en{EBITDA} הם \RRebitdaMarginFinal{} אחוזים, כלומר \RRebitdaFinal{} מיליון \shekel{}. הפער בין השניים הוא עלות המטה. היא גדלה לאט יותר מההכנסה, וזה מקור המינוף התפעולי. התפתחות מבנה ההכנסות מוצגת באיור~\ref{fig:revenue-mix}.

\begin{figure}[htbp]
\centering
\includegraphics[width=0.92\linewidth]{fig_revenue_mix}
\caption{הרכב ההכנסות לפי שנים: תוצרת טרייה מול רישוי \en{SaaS}. חלקו של ה-\en{SaaS} גדל אחרי 2029.}
\label{fig:revenue-mix}
\end{figure}

\section{נקודת איזון}

העסק נכנס לרווחיות תפעולית כשהתרומה הכוללת של הגגות מכסה את עלות המטה. אם $O$ היא עלות המטה השנתית ו-$c$ התרומה השנתית של גג בוגר, מספר הגגות המינימלי הוא:

\begin{equation}\label{eq:break-even}
N^{*} = \frac{O}{c}
\end{equation}

כאשר:
\begin{itemize}
  \item $N^{*}$ הוא מספר הגגות הבוגרים הנדרש לאיזון;
  \item $O$ היא עלות המטה וההנהלה של אותה שנה;
  \item $c$ היא התרומה של גג בוגר, \RRroofContribution{} אלף \shekel{} בשנה, אחרי דמי גג ועלות תפעול.
\end{itemize}

משוואה~\eqref{eq:break-even} מבהירה את הרגישות של העסק. כל עלייה בתרומה לגג מורידה את מספר הגגות הדרוש, וכל תוספת למטה מעלה אותו. לפי המודל, \en{EBITDA} הופך לחיובי ב-\RRbreakEvenYear{}, כשהצי גדול מספיק כדי לכסות את המטה, והוא עדיין לא ברמה שהמטה צריך לה בשנת ההקמה.

\section{תזרים מזומנים ועקומת \texorpdfstring{\en{J}}{J}}

באיור~\ref{fig:cash-curve} רואים את העקומה האופיינית לעסק עתיר השקעה. בשנים הראשונות כל גג חדש דורש השקעה לפני שהוא מניב, והמטה כבר פועל במלוא העלות. המזומן יורד, מגיע לשפל, ורק אחר כך מתחיל לעלות. הצורה הזאת נקראת עקומת \en{J}.

\begin{figure}[htbp]
\centering
\includegraphics[width=0.92\linewidth]{fig_cash_curve}
\caption{יתרת המזומן בסוף כל שנה, בצירוף מקורות ההון. היתרה הנמוכה ביותר היא \RRminCash{} מיליון \shekel{} ב-\RRminCashYear{}.}
\label{fig:cash-curve}
\end{figure}

יתרת המזומן הנמוכה ביותר היא \RRminCash{} מיליון \shekel{} ב-\RRminCashYear{}, והיא נשארת חיובית רק בזכות סבב \en{Seed} של \RRseed{} מיליון \shekel{} שמגויס עכשיו. סבב \en{Series A} של \RRseriesA{} מיליון \shekel{} ב-2029 מממן את שלב ההרחבה, וסך ההון שהחברה מגייסת הוא \RRequityTotal{} מיליון \shekel{}.

\begin{insightbox}[המשמעות למשקיע]
השפל של \RRminCash{} מיליון \shekel{} הוא מרווח ביטחון קטן. כל עיכוב בהתקנות או ירידה במחיר מקרבים את המזומן לאפס לפני הסבב הבא. לכן המימון בסבב \en{Series A} מועבר בתשלומים לפי אבני דרך (פרק~\ref{ch:risks}).
\end{insightbox}

\section{שווי נוכחי נקי}

כדי לתמצת את התזרים לערך אחד, המודל מחשב שווי נוכחי נקי של תזרימי המזומנים החופשיים, ועליהם שווי יציאה בסוף התקופה:

\begin{equation}\label{eq:npv}
\mathrm{NPV} = \sum_{t=1}^{T}\frac{\mathrm{FCF}_t}{(1+r)^{t}} + \frac{\mathrm{EBITDA}_T \cdot M}{(1+r)^{T}}
\end{equation}

כאשר:
\begin{itemize}
  \item $\mathrm{FCF}_t$ הוא התזרים החופשי בשנה $t$, אחרי מס והשקעות;
  \item $r$ הוא שיעור ההיוון, \RRdiscount{} אחוזים;
  \item $M$ הוא מכפיל היציאה על \en{EBITDA}, \RRexitMultiple;
  \item $T$ הוא אורך התקופה, חמש שנים.
\end{itemize}

לפי משוואה~\eqref{eq:npv}, התוצאה היא שווי נוכחי של \RRnpv{} מיליון \shekel{}. חשוב להבין שחלק גדול ממנו מגיע משווי היציאה ולא מהתזרים עצמו, כי בשנים הראשונות התזרים החופשי שלילי בגלל ההשקעות. לכן ה-\en{NPV} רגיש במיוחד ל-\en{EBITDA} של השנה האחרונה ולמכפיל.

\section{ניתוח רגישות}

כדי לבדוק איזה משתנה משפיע יותר מכל, מזיזים כל הנחה בעשרים אחוזים לכל כיוון ובודקים את השינוי ב-\en{NPV}. התוצאה באיור~\ref{fig:tornado}.

\begin{figure}[htbp]
\centering
\includegraphics[width=0.92\linewidth]{fig_tornado}
\caption{תרשים טורנדו: השפעת שינוי של עשרים אחוזים בכל הנחה על ה-\en{NPV}.}
\label{fig:tornado}
\end{figure}

המסקנה ברורה: ההשפעה הגדולה ביותר היא של \RRtornadoFirst{}, כי הם משפיעים על כל גג בכל שנה. אחריהם בא הגורם הבא בדירוג, \RRtornadoSecond{}, ואז \RRtornadoThird{}. ההשפעה הקטנה ביותר היא של \RRtornadoLast{}.

\begin{riskbox}[מגבלות המודל]
המודל הוא אוסף של הנחות. תפוקה של \RRyield{} ק"ג למ"ר ומחיר של \RRprice~\shekel{} לק"ג עדיין לא אומתו במערך מסחרי בקנה מידה. ההנחה שגג חדש פועל בחצי קיבולת בשנה הראשונה ושמכפיל היציאה הוא \RRexitMultiple{} מבוססת על הערכת הצוות. המספרים משמשים לכיול ההחלטות ולא כהתחייבות.
\end{riskbox}

\section{סיכום}

הכלכלה ברמת הגג חיובית, עם החזר השקעה של \RRroofPayback{} שנים. ברמת החברה, התוצאה תלויה בקצב פריסה ובשמירה על מחיר ותפוקה. המשקיע צריך לעקוב אחרי שלושה מדדים: תפוקה למ"ר, מחיר ממוצע לק"ג ומספר גגות שנחתמו. כל אחד מהם מתעדכן חודשית (ראו פרק~\ref{ch:team}), ויחד הם קובעים אם התוכנית נשארת על המסלול.
